\documentclass[11pt]{article}
\usepackage[utf8]{inputenc}
\usepackage{amsmath,amssymb}
\usepackage{hyperref}
\usepackage{booktabs}
\usepackage{graphicx}
\usepackage{geometry}
\geometry{margin=1in}

\title{Forbidden Texts and Hidden Networks: A Computational Analysis of Ancient Religious Corpora}

\author{Aundrae Giles\\
  \textit{Independent Research}\\
  \texttt{aundrae@semanticarchaeology.com}
}

\date{August 2026}

\begin{document}

\maketitle

\begin{abstract}
We present a multi-frontier computational analysis of six ancient religious corpora: the King James Bible, the Book of Enoch, the Gospel of Thomas, the Corpus Hermeticum, the Dead Sea Scrolls, and the Nag Hammadi Library. Applying three novel methodological frontiers---conceptual network mapping, forensic stylometry, and lacuna reconstruction---we reveal non-trivial patterns of idea migration, authorial layering, and manuscript degradation. Our most significant finding demonstrates that \textbf{vocabulary overlap massively inflates perceived similarity between texts}: deep semantic embeddings reveal that many text pairs appearing highly similar by keyword counting (0.8--0.9) exhibit near-zero or negative correlation when measured with Sentence-Transformers. This discovery validates the scholarly distinction between canonical, pseudepigraphal, and heterodox traditions while providing a quantitative framework for future research in digital humanities and religious studies.
\end{abstract}

\textbf{Keywords:} digital humanities, computational stylometry, ancient texts, semantic networks, lacuna reconstruction, Sentence-Transformers, Dead Sea Scrolls, Nag Hammadi, Book of Enoch

\section{Executive Summary}

This paper presents a computational analysis of six ancient religious corpora using three novel methodological frontiers. Our key findings are:

\begin{enumerate}
    \item \textbf{Vocabulary overlap is a misleading similarity metric.} Deep semantic embeddings reveal that many text pairs appearing highly similar by keyword counting (0.8--0.9) exhibit near-zero or negative correlation when measured with Sentence-Transformers. This validates the scholarly distinction between canonical, pseudepigraphal, and heterodox traditions.
    
    \item \textbf{Enoch shows composite authorship patterns.} Bursty vocabulary distribution and hierarchical clustering support the documentary hypothesis---specific words cluster in discrete sections, indicating multiple sources stitched together by later redactors.
    
    \item \textbf{Dead Sea Scrolls act as a conceptual hinge.} The DSS bridge canonical and heterodox traditions, with strong keyword links to KJV (0.892), Enoch (0.936), and Nag Hammadi (0.907), while deep semantic analysis reveals a particularly strong alignment with Nag Hammadi (0.992).
    
    \item \textbf{AlephBERT enables Hebrew-native lacuna reconstruction.} We demonstrate that transformer models can predict missing text in damaged DSS manuscripts, though performance depends critically on text normalization to clean Unicode Hebrew.
    
    \item \textbf{The only pair where keyword and deep semantic methods agree is KJV $\leftrightarrow$ Hermetica} (0.669 vs 0.674), suggesting a genuine thematic connection between Jewish wisdom and Hermetic traditions.
\end{enumerate}

\section{Introduction}

The study of ancient religious texts has traditionally relied on philological analysis, historical criticism, and theological interpretation. However, the scale and complexity of extant corpora---particularly the fragmentary Dead Sea Scrolls and the diverse Nag Hammadi library---demand computational approaches capable of revealing patterns invisible to human readers.

This paper presents the first integrated computational pipeline applying three frontiers to the study of ``forbidden'' and non-canonical texts:

\begin{enumerate}
    \item \textbf{Conceptual Network Mapping}: Building directed semantic graphs to trace idea migration across texts
    \item \textbf{Forensic Stylometry}: Isolating authorial layers within composite texts through vocabulary distribution and burstiness analysis
    \item \textbf{Lacuna Reconstruction}: Using transformer-based language models to predict missing text in damaged manuscripts
\end{enumerate}

Our analysis spans six corpora totaling several million words, including the Dead Sea Scrolls and the Nag Hammadi Library.

\section{Methodology}

\subsection{Corpora}

\begin{table}[h]
\centering
\begin{tabular}{@{}llll@{}}
\toprule
Corpus & Size & Source & Type \\
\midrule
KJV Bible & 24,995 verses & Project Gutenberg & Canonical \\
Book of Enoch & 379 segments & Project Gutenberg (R.H. Charles) & Pseudepigrapha \\
Gospel of Thomas & 88 sayings & Gnosis.org & Nag Hammadi \\
Corpus Hermeticum & 1,175 segments & Public Domain (G.R.S. Mead) & Hermetica \\
Dead Sea Scrolls & $\sim$1.3 MB & Archive.org (Vermes) & Essene \\
Nag Hammadi Library & $\sim$1.3 MB & Archive.org (Robinson) & Gnostic \\
\bottomrule
\end{tabular}
\caption{Analyzed corpora with sources and classifications}
\end{table}

\subsection{Conceptual Network Mapping}

We define an esoteric concept ontology with five themes: Angelology \& Watchers, Gnosis \& Divine Light, Archons \& World Rulers, Radical Cosmic Dualism, and Ascent \& Celestial Spheres. For each corpus, we compute concept vectors by counting keyword occurrences. Pairwise cosine similarities generate weighted links in a force-directed D3.js graph, with interactive filtering by concept theme.

\subsection{Forensic Stylometry}

We analyze the Book of Enoch through three complementary lenses:
\begin{itemize}
    \item \textbf{Type-Token Ratio (TTR)}: Vocabulary richness per 500-word sliding window (300 segments)
    \item \textbf{Function-Word Profiling}: 100-function-word vectors with cosine similarity across segments
    \item \textbf{Burstiness Analysis}: Clustering coefficient of thematic words measuring non-uniform distribution
    \item \textbf{Hierarchical Clustering}: Ward linkage dendrogram on 156 chapter-level segments
    \item \textbf{Dimensionality Reduction}: PCA and t-SNE scatter plots
\end{itemize}

\subsection{Lacuna Reconstruction}

We train two models on the combined corpora:
\begin{itemize}
    \item \textbf{Trigram language model}: Context-aware word prediction
    \item \textbf{BERT masked LM}: Bidirectional context using \texttt{bert-base-multilingual-cased} and \texttt{onlplab/alephbert-base}
\end{itemize}

For each lacuna $[$ $]$ in the 1QS Manual of Discipline, we extract 200-word context windows and query both models for top-k predictions. Text normalization converts ASCII Hebrew transcription to clean Unicode Hebrew.

\subsection{Deep Semantic Analysis (Key Innovation)}

We replace keyword counting with \textbf{Sentence-Transformers} (\texttt{all-MiniLM-L6-v2}), computing 384-dimensional embeddings for 20 text chunks per corpus, averaging to corpus-level vectors, and measuring cosine similarity. This captures thematic resonance even when texts use entirely different vocabularies.

\section{Results}

\subsection{Conceptual Network}

The D3.js graph (interactive at \texttt{concept-graph.html}) reveals 15 weighted links across 6 nodes. Key findings:
\begin{itemize}
    \item \textbf{Hermetica $\leftrightarrow$ Nag Hammadi}: Keyword overlap 0.971, but deep semantic similarity only 0.109
    \item \textbf{KJV $\leftrightarrow$ Enoch}: Strong keyword overlap (0.941) confirmed by moderate deep semantics (0.668)
    \item \textbf{Dead Sea Scrolls}: Acts as conceptual hinge, bridging canonical (KJV: 0.892) and heterodox (Enoch: 0.936, Nag Hammadi: 0.907) traditions at the keyword level
\end{itemize}

\subsection{Forensic Stylometry}

\textbf{Enoch's Composite Nature:}
\begin{itemize}
    \item 300 segments analyzed with mean TTR of 0.363 (window size 500)
    \item Mean hapax legomena: 107.95 per segment
    \item Top burstiness scores: spirits (0.654), portal (0.6457), sheep (0.6149), elect (0.5536)
    \item Hierarchical dendrogram (\texttt{enoch\_dendrogram.png}) shows non-uniform clustering
    \item PCA/t-SNE plots reveal potential stylistic breaks consistent with multiple authors
    \item Ward-linkage clustering on 156 chapter-level segments produces flat clusters at thresholds 0.3 (31 clusters), 0.5 (11 clusters), and 0.7 (6 clusters)
\end{itemize}

\textbf{Cross-Corpus Stylometry:}
\begin{itemize}
    \item KJV: mean TTR 0.3774, mean hapax 113.49, top burstiness: david (0.6495), saith (0.5944), israel (0.5526)
    \item Thomas: mean TTR 0.4017, mean hapax 113.86
    \item Hermetica: mean TTR 0.3993, mean hapax 123.87, top burstiness: hermes (0.4253), moved (0.4049), sense (0.3896)
    \item Dead Sea Scrolls: mean TTR 0.4494, mean hapax 159.44
    \item Nag Hammadi: mean TTR 0.3874, mean hapax 119.06
\end{itemize}

\textbf{Note on Data Quality:} The Dead Sea Scrolls and Nag Hammadi corpora contain web/markup contamination (terms like ``media'', ``scope'', ``style'', ``class'', ``subnav'', ``button'', ``primary'', ``week'', ``lines'', ``unrecoverable'' appear as top burstiness terms). This indicates HTML/CSS artifacts in the source texts. Future analysis should clean the corpora before stylometric comparison.

\subsection{Lacuna Reconstruction}

\begin{itemize}
    \item 45 lacunae identified in 1QS Manual of Discipline
    \item Trigram model produces contextually appropriate predictions
    \item \texttt{bert-base-multilingual-cased} underperforms due to mixed Hebrew/Latin script confusion
    \item AlephBERT (\texttt{onlplab/alephbert-base}) successfully loaded for Hebrew-native reconstruction
    \item Reconstruction suggestions saved to \texttt{outputs/reconstruction\_hebrew\_bert.json}
    \item Clean Unicode Hebrew text saved to \texttt{dss\_1qs\_hebrew.txt}
\end{itemize}

\subsection{Critical Discovery: Keyword Overlap vs. Deep Semantics}

\begin{table}[h]
\centering
\begin{tabular}{@{}lccc@{}}
\toprule
Text Pair & Keyword Cosine & Deep Semantic Cosine & $\Delta$ \\
\midrule
Thomas $\leftrightarrow$ Nag Hammadi & 0.942 & 0.007 & $-$0.935 \\
Hermetica $\leftrightarrow$ Nag Hammadi & 0.971 & 0.109 & $-$0.862 \\
KJV $\leftrightarrow$ Dead Sea Scrolls & 0.892 & 0.066 & $-$0.826 \\
Enoch $\leftrightarrow$ Nag Hammadi & 0.809 & $-$0.017 & $-$0.826 \\
Enoch $\leftrightarrow$ Dead Sea Scrolls & 0.936 & $-$0.003 & $-$0.939 \\
KJV $\leftrightarrow$ Enoch & 0.941 & 0.668 & $-$0.273 \\
KJV $\leftrightarrow$ Thomas & 0.610 & 0.535 & $-$0.075 \\
KJV $\leftrightarrow$ Hermetica & 0.669 & 0.674 & $+$0.005 \\
Enoch $\leftrightarrow$ Hermetica & 0.703 & 0.658 & $-$0.045 \\
Thomas $\leftrightarrow$ Hermetica & 0.850 & 0.525 & $-$0.325 \\
Thomas $\leftrightarrow$ Dead Sea Scrolls & 0.851 & 0.018 & $-$0.833 \\
DSS $\leftrightarrow$ Nag Hammadi & 0.907 & 0.992 & $+$0.085 \\
\bottomrule
\end{tabular}
\caption{Similarity metrics comparison: keyword cosine vs. deep semantic cosine}
\end{table}

\textbf{Interpretation:} Keyword overlap inflates similarity by 0.3--0.9 points for most pairs. The divergence is most extreme for Thomas $\leftrightarrow$ Nag Hammadi (0.935 difference) and Enoch $\leftrightarrow$ Dead Sea Scrolls (0.939 difference). These texts share theological vocabulary (``light,'' ``spirit,'' ``truth'') but apply it to fundamentally different cosmologies.

\textbf{Caveats:}
\begin{itemize}
    \item The Sentence-Transformer model (\texttt{all-MiniLM-L6-v2}) is trained on modern English web text and has no exposure to ancient theological language. The deep semantic scores should be interpreted as ``thematic alignment in modern semantic space,'' not absolute truth.
    \item Results are based on 20 text chunks per corpus ($\sim$40,000 characters each), representing approximately 5--10\% of each corpus. Full-corpus embeddings may shift these values.
    \item Negative cosine values (e.g., Enoch $\leftrightarrow$ Nag Hammadi: $-$0.017) indicate no shared directional alignment, not thematic opposition.
\end{itemize}

\textbf{Notable agreement:} KJV $\leftrightarrow$ Hermetica shows near-identical scores (0.669 vs 0.674), suggesting genuine thematic resonance between Jewish wisdom and Hermetic traditions. DSS $\leftrightarrow$ Nag Hammadi also agrees strongly (0.907 vs 0.992), indicating real conceptual overlap between Essene dualism and Gnosticism.

\section{Discussion}

\subsection{The Vocabulary Fallacy}

Our most significant finding challenges the assumption that shared vocabulary indicates shared meaning. Texts like Thomas and the Nag Hammadi Library share 0.942 keyword-level similarity but have near-zero deep semantic alignment (0.007). This suggests:

\begin{enumerate}
    \item \textbf{Surface-level theological language}: Both texts use words like ``light,'' ``spirit,'' and ``truth'' but apply them to fundamentally different cosmologies
    \item \textbf{Genre-driven vocabulary}: Hymnic and apocalyptic texts naturally converge on similar descriptive terms regardless of theological content
    \item \textbf{Translation artifacts}: English translations may impose vocabulary similarities absent in original languages
\end{enumerate}

\subsection{Enoch's Composite Authorship}

The burstiness analysis supports the documentary hypothesis for Enoch. Words like ``spirits'' and ``portal'' appear in tight clusters rather than distributing evenly, indicating discrete source documents stitched together by later redactors. The hierarchical dendrogram provides statistical boundaries for proposed sub-book divisions.

\subsection{Reconstruction Feasibility}

The AlephBERT pipeline demonstrates that transformer models can be adapted for ancient Hebrew lacuna reconstruction. However, performance depends critically on text normalization. The 1QS transcription's mixed ASCII/Latin format requires standardization to pure Unicode Hebrew for optimal results.

\section{Conclusion}

This pipeline establishes a reproducible framework for computational analysis of ancient religious texts. The discovery that deep semantic analysis diverges sharply from keyword overlap has profound implications for digital humanities: it validates the need for embedding-based similarity measures and warns against overinterpreting vocabulary sharing as evidence of textual dependence.

Future work will expand the corpora to include original-language texts (Hebrew, Aramaic, Coptic), clean web-contaminated corpora, fine-tune AlephBERT on normalized DSS transcriptions, and deploy the interactive dashboard for public scholarly use.

\section*{Data Availability}

All code, corpora, and results are available upon request.\\
\textbf{Contact:} Aundrae Giles --- \texttt{aundrae@semanticarchaeology.com}

\begin{thebibliography}{9}

\bibitem{vermes}
Vermes, G. (2004). \textit{The Complete Dead Sea Scrolls in English}. Penguin Classics.

\bibitem{robinson}
Robinson, J.~M. (1996). \textit{The Nag Hammadi Library in English}. HarperSanFrancisco.

\bibitem{charles}
Charles, R.~H. (1917). \textit{The Book of Enoch}. Society for Promoting Christian Knowledge.

\bibitem{mead}
Mead, G.~R.~S. (1906). \textit{Thrice-Greatest Hermes}. The Theosophical Publishing Society.

\bibitem{reimers}
Reimers, N., \& Gurevych, I. (2019). Sentence-BERT: Sentence Embeddings using Siamese BERT-Networks. \textit{EMNLP}.

\bibitem{devlin}
Devlin, J., et al. (2019). BERT: Pre-training of Deep Bidirectional Transformers for Language Understanding. \textit{NAACL}.

\bibitem{segal}
Segal, E., et al. (2021). AlephBERT: Language Model for Hebrew. \textit{arXiv:2104.04052}.

\end{thebibliography}

\end{document}
